\subsection{Six triangular seven-element sources with actual state quotients}

\begin{theorem}[Proved: the seven-element triangular quotient family]
\label{CP2:A-seven-triangular-family-right-zero}
Let $P=\{1,e_{00},e_{01},e_{10},e_{11}\}$ be the
identity-adjoined rectangular band, put $K=P\setminus\{1\}$,
and abbreviate $e_0=e_{00}$ and $e_1=e_{01}$.
For each row of the following table, form $M=P\amalg\{a,b\}$.
The two columns headed $\pi$ specify the mixed products by
\[
 ak=\pi(a)k,\quad ka=k\pi(a),\qquad
 bk=\pi(b)k,\quad kb=k\pi(b)\qquad(k\in K);
\]
the remaining four columns specify all products of $a,b$.
The first column gives the index in the frozen positive
order-seven catalog.
\[
\begin{array}{c|cc|cccc}
\text{index}&\pi(a)&\pi(b)&a^2&ab&ba&b^2\\\hline
3&e_0&e_0&e_0&e_0&a&b\\
15&e_0&e_0&a&b&a&b\\
19&e_0&1&e_0&a&e_0&b\\
28&1&1&a&a&b&b\\
34&e_0&e_1&e_0&e_1&a&b\\
39&e_0&e_1&a&b&a&b
\end{array}
\]
Let $N=R_3$ be the identity-adjoined two-element right-zero
band.  For every finite nonempty left-$N$, right-$M$ biset
$S$, bounded rational homology of its collage implies that
all rational degree-two products vanish.  There is no
restriction on $|S|$ or on the second Betti number.
\end{theorem}
\begin{proof}
The displayed assignment, together with $\pi|_P=\operatorname{id}$,
is a unital monoid retraction $\pi:M\to P$ in every row.
The mixed products respect it by definition, and the four
remaining products respect it directly from the table.
We shall construct a corresponding actual quotient of the
connecting state set and compute the relative homology of
the actual collage functor.

Put $R=\mathbb QM$, $B=\mathbb QP$, and
$\Delta=B/\mathbb QK$, regarded as a right $B$ module.
Use the following elements of $R$.  In each row $q,r$ are
monoid idempotents, $qr=rq=r$, and $h=q-r$; put $p=1-h$.
The last two columns give the scalars by which left
multiplication by $a,b$ acts on each of the one-dimensional
spaces $\mathbb Qh$ and $\mathbb Qz$.
\[
\begin{array}{c|cc|c|cc}
\text{index}&q&r&z&a&b\\\hline
3&b&e_0&a-e_0&0&1\\
15&a&e_0&b-a&1&1\\
19&1&b&a-e_0&0&0\\
28&1&a&b-a&0&0\\
34&b&e_1&a-e_0&0&1\\
39&a&e_0&b-a-e_1+e_0&1&1
\end{array}
\]
Multiplication using the defining table gives, in every row,
\[
\begin{gathered}
 h^2=h,\qquad z^2=0,\qquad hz=z,\qquad zh=0,\\
 hK=Kh=zK=Kz=0,\qquad
 Rh=\mathbb Qh,\qquad
 Rp=\mathbb Q\{p,e_{00},e_{01},e_{10},e_{11},z\}.
\end{gathered}
\]
The seven elements $p,h,z,e_{00},e_{01},e_{10},e_{11}$
are a basis of $R$.  Since $pz=0$ and $zp=z$, the corner
$pRp$ has basis $p,e_{00},e_{01},e_{10},e_{11}$ and is
identified by $\pi$ with $B$, with $p$ as its identity.
The identical left scalars on $h,z$, and their annihilation
by $K$, give the exact left-$R$, right-$B$ bimodule sequence
\[
 0\longrightarrow Rh\otimes_{\mathbb Q}\Delta
 \xrightarrow{\ h\otimes1\mapsto z\ }Rp
 \xrightarrow{\pi}B\longrightarrow0.                 \tag{1}
\]
Here the right $B$ action on $Rp$ is through $pRp$;
$zK=0$ and $zp=z$ verify right linearity.  Both terms
preceding $B$ are projective as left $R$ modules.
The displayed basis and scalar calculations prove (1)
for all six rows, not just for their contracted algebras.

We next identify its state tensor by actual sets.  Write
$A(s)=sa$, $B_s(s)=sb$, $E_i(s)=se_i$, $Q(s)=sq$ and
$R_s(s)=sr$; the subscript distinguishes these functions
from the two algebras.  Form the quotient set $T$ by
identifying $t$ with $R_s(t)$ for every $t\in Q(S)$.
The relations $Q^2=Q$, $R_s^2=R_s$ and
$QR_s=R_sQ=R_s$ give canonical representatives
\[
 T=(S\setminus Q(S))\amalg R_s(S).
\]
Write $\rho:S\to T$ for the quotient.  It is $N$
equivariant because all the functions commute with $N$.
It is also a right $P$ quotient: for $t=sq$ and $k\in K$,
\[
 tk=sqk=srk=(R_s(t))k,
\]
using $hK=0$.  The identity of $P$ acts as the identity.
Thus $T$ is an actual left-$N$, right-$P$ biset.

For $L=\mathbb Q[S]$, projection to $\mathbb Q[T]$
restricts to a natural isomorphism
\[
 Lp\xrightarrow{\cong}\mathbb Q[T].                    \tag{2}
\]
An integral basis on the left is
\[
 [s]-[Q(s)]+[R_s(s)]\quad(s\notin Q(S)),
 \qquad[c]\quad(c\in R_s(S)).
\]
These vectors project to the distinct indicated representatives
of $T$ and span $Lp$ by applying $p$ to every state.
The isomorphism is left-$N$ linear and right-$B$ linear:
the right action of $p$ is the identity, and multiplication
by every band element is respected by $\rho$ and $pK=K$.
Likewise $Lh$ has basis
\[
 v_t=[t]-[R_s(t)]\qquad(t\in Q(S)\setminus R_s(S)).
\]

Tensoring (1) gives the natural right-$B$ complex
\[
 F_L=[\,Lh\otimes_{\mathbb Q}\Delta
                   \xrightarrow{\ v\mapsto vz\ }Lp\,]
\]
in degrees one and zero.  Under (2), every basis vector
$v_t$ maps to the difference of two actual vertices of $T$.
The complete list is
\[
\begin{array}{c|c|c}
\text{index}&\rho(v_tz)&\text{equivalent relations for all }s\in S\\\hline
3&[\rho(A(t))]-[\rho(E_0(t))]&\rho(A(s))\sim\rho(E_0(s))\\
15&[\rho(B_s(t))]-[\rho(E_0(t))]&\rho(B_s(s))\sim\rho(E_0(s))\\
19&[\rho(A(t))]-[\rho(E_0(t))]&\rho(A(s))\sim\rho(E_0(s))\\
28&[\rho(B_s(t))]-[\rho(A(t))]&\rho(B_s(s))\sim\rho(A(s))\\
34&[\rho(A(t))]-[\rho(E_0(t))]&\rho(A(s))\sim\rho(E_0(s))\\
39&[\rho(B_s(t))]-[\rho(E_1(t))]&\rho(B_s(s))\sim\rho(E_1(s))
\end{array}                                                    \tag{3}
\]
Here $\rho$ on a linear combination denotes its linear extension.
To check that the last column gives exactly the same relations,
in rows $3,34$ take $t=B_s(s)$: the table gives
$AB_s=A$ and $E_0B_s=E_0$.  In rows $15,39$ take
$t=A(s)$: it gives $B_sA=B_s$ and $E_iA=E_i$ for
the indicated $i$.  In rows $19,28$, $Q$ is the identity,
so the vectors already range over every state outside $R_s(S)$.
In all six cases the corresponding difference is zero on
$R_s(S)$, so no relation is lost by omitting that subset.
These identities also verify each entry of the middle column.
For row $39$, the term $-a+e_0$ disappears after projection
to $T$, and the two terms obtained by applying $b,e_1$
to $E_0(t)$ cancel; this leaves precisely its stated difference.

Let $U$ be the set of connected components of the graph on
$T$ having the pairs in (3) as its edges.  Equivalently,
impose exactly the relations in the last column.  The
quotient is $N$ equivariant.  It is preserved by $P$ because
the two functions in each relation have the same composite
with every right band actor; this follows from the two
corresponding mixed-product columns.  Thus $U$ is an actual
left-$N$, right-$P$ biset.  The first quotient $S\to T$
already identifies the actions of $q,r$, and (3) imposes
the remaining relation between $a,b$ and their images
under $\pi$.  Consequently $q_S:S\to U$ satisfies
\[
 q_S(sm)=q_S(s)\pi(m)\qquad(m\in M).
\]
It defines an actual two-object functor
\[
 q:(M,N,S)\longrightarrow(P,N,U).                       \tag{4}
\]
No linear quotient is substituted for a set quotient here.

The cokernel of the degree-one map of $F_L$ is exactly
$\mathbb Q[U]$: the span of edge differences in a graph
is the kernel of the free-vector-space map to its components.
Let $W$ be the kernel of that map on the coefficient space
$Lh$.  It is a finite-dimensional left $\mathbb QN$ module.
All of the degree-one right $B$ action factors through
$\Delta$, so the kernel of the natural chain surjection
\[
 F_L\longrightarrow\mathbb Q[U]
\]
is naturally quasi-isomorphic to
$(W\otimes_{\mathbb Q}\Delta)[1]$.  This remains true
when the graph in (3) has loops, repeated edges or cycles;
there is no injectivity assumption on its incidence map.

The remainder of the comparison can be made in the same
explicit degrees as in
Theorem~\ref{CP2:A-seven-source-three-right-zero}.
The length-two $P_5$ projective resolution has terms
$Be$, $B(1-e)$, $Be$ in degrees zero, one, two.
Tensoring with $\Delta$ kills the two outer terms and
leaves one copy of $\mathbb Q$ in degree one.  Hence
\[
 \Delta\otimes_B^{\mathbf L}\mathbb Q\simeq\mathbb Q[1].
\]
Derived associativity applied to (1), naturally with respect
to the actual quotient maps, now identifies the fiber of
the rational chain map of bridges
\[
 B(N,S,M)\longrightarrow B(N,U,P)
\]
with $(\mathbb Q\otimes_{\mathbb QN}^{\mathbf L}W)[2]$.
Its mapping-cone group in degree $j$ is therefore
$\operatorname{Tor}^{\mathbb QN}_{j-3}(\mathbb Q,W)$.
Taking $L=\mathbb Q$ in (1) gives $Lh=0$ and $Lp=\mathbb Q$;
thus the actual corner map $BM\to BP$ is a rational
homology equivalence.

In the chain mapping-cylinder models for the two collages,
the bridge terms enter one degree above the corner terms.
The cone of the corner comparison is acyclic.  The long
exact sequence of the resulting map of mapping cones gives
\[
 \begin{aligned}
 \cdots&\longrightarrow H_j(B(M,N,S);\mathbb Q)
 \longrightarrow H_j(B(P,N,U);\mathbb Q)\\
 &\longrightarrow\operatorname{Tor}^{\mathbb QN}_{j-4}(\mathbb Q,W)
 \longrightarrow H_{j-1}(B(M,N,S);\mathbb Q)
 \longrightarrow\cdots .
 \end{aligned}                                           \tag{5}
\]
In particular (4) induces an isomorphism on rational
$H_2$ and hence a natural isomorphism on rational $H^2$.
The trivial right $\mathbb QN$ module has projective
dimension at most one, by the explicit $R_3$ resolution
in the cited theorem.  Thus the relative groups in (5)
can occur only in degrees four and five.  Boundedness of
the original collage transfers to $(P,N,U)$.
Corollary~\ref{CP2:A-baseline-right-zero-three-six} kills
all rational degree-two products of the latter collage.
Every second class of the original is pulled back along
the actual functor (4), so all its products vanish too.

The separate literal binding reconstructs every product
of all six tables, the quotient maps and every displayed
algebra identity in
\path{converse/scope18/seven_triangular_family_literal_binding.json}.
It uses identity $0$, band labels $1,2,3,4$, and $a=5,b=6$.
This binds the stated catalog indices without using any
classification of the connecting actions.
\end{proof}
